%%%%%
%
\section{Sous-�chantillonnage r�gulier (param�tre \texttt{double})}

\itemstyle{Op\'erateur~:} \operfmt{ Ech\_dbl}
\index{Ech\_dbl}

%%%%%OBJET%%%%
\itemstyle{Objet~:}
Echantillonnage d'un param\`etre \texttt{double} \`a une
fr\'equence inf\'erieure. La fr\'equence 
finale  doit \^etre un diviseur
de la fr\'equence d'origine.

%%%%%ARGUMENTS%%%%
\bigskip
%%%%%ENTREES%%%%
\begin{tabarg}
\multicolumn{3}{l}{\titretab {Entr\'ees}}\\
\hline
$X$ & Vecteur, \texttt{double} & Param\`etre d'entr\'ee\\
$R_f$ & Coeff., \texttt{float} & Rapport des fr\'equences $ F_{sortie}\over
F_{entree}$, dans l'intervalle $]0, 1]$ \\
$Decal$ & Coeff., \texttt{float} & D\'ecalage (\textit{nombre de points}), compris entre 0 et $1/R_f$\\
\end{tabarg}

\medskip
%%%%%SORTIES%%%%
\begin{tabarg}
\multicolumn{3}{l}{\titretab {Sorties}}\\
\hline
$Xe$ & Vecteur, \texttt{double} & Param\`etre \'echantillonn\'e (m\^eme unit\'e que $X$)\\
\end{tabarg}


%%%%%FORMULE%%%%
\itemstyle{Formule}

%%Ech
$$
Xe[i] = X[i/R_f + Decal]  \qquad i=0, Nbre\_echant(X)*R_f -1.
$$


%%%%%COMMENTAIRE%%%%
\itemstyle{Commentaire~:}

$R_f$ doit \^etre l'inverse d'un entier. 

Le nombre d'\'echantillons du vecteur de sortie est le produit du
nombre d'\'echantillons du vecteur d'entr\'ee par $R_f$, arrondi 
\`a l'entier sup\'erieur.

Le ref\_time du vecteur r�sultat ($Xe$) s'appelle time\_\$\{freq\}

%%%%%AUTEUR%%%%
\itemstyle{Auteur~:}
CNRM/GMEI/TRAMM

%%%%%REFERENCE%%%%
\itemstyle{R\'ef\'erences~:}
